\chapter{控制屏障函数与安全控制：CLF/CBF/QP}
\label{ch:clf-cbf}

% ════════════════════════════════════════════════════════════════
%  P4 控制部 · C1b 第 F 章 · 归卷四《控制理论基础》
%  主源：Tutorial 01_数学/40_控制理论/80_CLF_CBF与QP安全控制.md（2024 行）
%  设计稿：docs/控制部设计_C1b_F.md（§4 大纲 / §2 勘误 / §3 记号 / §5 cref / §6 bib）
%  主线：集合不变性 → CLF 引子 → CBF 主体 → 高相对度 → 统一 QP → 可行性 → 前沿 → 应用
%  与 D 章（Lyapunov/ISS/比较引理）严格去重：CBF 写成"安全对偶"，比较引理一般形式 cref D
% ════════════════════════════════════════════════════════════════

\noindent
本章把"安全"——系统永不离开容许集——锻造成与稳定性同等严格的控制证书。我们从集合不变性的 Nagumo 定理出发，引入控制 Lyapunov 函数（CLF）作为对偶引子，随后把控制屏障函数（CBF）建立为本书的核心安全工具，并将其化为可实时求解的二次规划（QP）。这条主线在卷四中承上启下：它是 \cref{ch:lyapunov-stability} 稳定性证书的安全对偶、\cref{ch:hj-reachability} 哈密顿-雅可比可达性的计算廉价近似，也是规划部安全滤波与强化学习运控部部署时安全层的理论母章。
\medskip

% ----------------------------------------------------------------
\section*{承上：从"会收敛"到"永远安全"}

\cref{ch:lyapunov-basics,ch:lyapunov-stability} 教会我们用 Lyapunov 函数 $V(x)$ 证明 \textbf{系统会收敛到目标}（$V(x)\to0$）。本章要教的是用控制屏障函数 $h(x)$ 证明 \textbf{系统永远不会离开安全区}（$h(x)\ge0$）。两者是一对精确对偶：前者刻画\emph{稳定性}（$t\to\infty$ 的最终行为），后者刻画\emph{安全性}（$\forall t\ge0$ 的全程行为）。

为什么必须把安全单独拎出来？\textbf{一个稳定的系统在收敛途中完全可能穿过危险区域}。机械臂从 A 点移到 B 点，Lyapunov 控制器保证它终将到达 B，但路径可能横穿工人作业区——稳定性只承诺终点，不承诺过程。打个比方：Lyapunov 函数像 GPS，告诉你"一定能到目的地"；CBF 像护栏，保证你"永远不会掉下悬崖"。GPS 可能给出一条贴着悬崖的最短路线，若无护栏，你虽终将抵达，过程中却可能坠落。CLF-CBF-QP 正是同时配备 GPS 与护栏的系统。

对做强化学习运动控制或具身智能的读者，这对双生工具尤为关键：RL 策略给出的"名义动作"未必安全，CBF-QP 作为安全滤波器（safety filter）能在毫秒级对其做\emph{最小修正}，从而保证硬约束被满足。

\begin{note}[本章定位：D 章的安全侧对偶]
本章与 \cref{ch:lyapunov-stability}（Lyapunov/LaSalle/ISS/反步）共享几乎相同的工具集合，只是把每个对象"翻转"到安全语义：$V\to h$、$\le\to\ge$、"收敛到 $0$"$\to$"保持 $\ge0$"、LaSalle 不变性 $\to$ Nagumo 不变性、ISS $\to$ ISSf。据此分工：类 $\mathcal K_\infty$ 函数定义与\textbf{比较引理的一般形式}（Khalil 引理 3.4/4.4）主证在 \cref{ch:lyapunov-basics}，本章只证其"安全特化"；ISS 主证在 \cref{ch:lyapunov-stability}，本章 ISSf（\cref{sec:cbf-robust}）引用之。读本章时，建议对每个 CBF 结论都自问："对应的 CLF/Lyapunov 版本是什么？"
\end{note}

\subsection*{CBF 与 CLF 的对偶总表}

\begin{table}[H]\centering\small
\caption{安全证书（CBF）与稳定证书（CLF）的逐项对偶}
\label{tab:cbf-clf-duality}
\begin{tabular}{lll}
\toprule
概念 & CLF（稳定性） & CBF（安全性） \\
\midrule
关心的集合 & 目标集 $\{V=0\}$ & 安全集 $\mathcal C=\{h\ge0\}$ \\
目标 & 吸引（$V\to0$） & 不变（$h\ge0$ 永成立） \\
核心不等式 & $\dot V\le-\gamma(V)$ & $\dot h\ge-\alpha(h)$ \\
不等式方向 & 上界（"至少以 $\gamma$ 下降"） & 下界（"至多以 $\alpha$ 下降"） \\
比较引理结论 & $V(t)\to0$ & $h(t)\ge0$ 保持 \\
工具定理 & LaSalle 不变性原理 & Nagumo 正不变性定理 \\
QP 角色 & CLF-QP（保收敛） & CBF-QP（保安全） \\
冲突时优先级 & 可松弛（软约束） & 不可松弛（硬约束） \\
\bottomrule
\end{tabular}
\end{table}

\begin{note}[本章记号约定]
全书控制仿射系统记 $\dot x=f(x)+g(x)u$，$x\in\mathbb R^n$，$u\in\mathcal U\subseteq\mathbb R^m$，$f,g$ 局部 Lipschitz；李导数 $L_fh=\nabla h\cdot f$、$L_gh=\nabla h\cdot g$（后者为 $1\times m$ 行向量）。安全函数记 $h$、CLF 记 $V$、倒数型屏障记 $B$，与 \cref{ch:lyapunov-basics} 的 $V$/类 $\mathcal K$ 记号对齐。为消歧，本章钉死四处符号：\emph{(i) 重力加速度记 $\mathsf g$}（区别于输入矩阵 $g(x)$，见四旋翼例 \cref{ex:cbf-quad}）；\emph{(ii) 衰减率 $\gamma$} 兼表 CLF 衰减率与 CBF 线性类 $\mathcal K$ 斜率（二者同质可共用），但 \cref{sec:cbf-unify} 的 CBVF 折扣率另记 $\gamma_{\mathrm{cbvf}}$；\emph{(iii) 位置用黑体 $\mathbf p$}、松弛惩罚权用细体标量 $p$；\emph{(iv) 相对度记 $r$}、几何半径统一记 $R$。
\end{note}

% ════════════════════════════════════════════════════════════════
\section{集合不变性与 Nagumo 定理}
\label{sec:cbf-nagumo}

\subsection{正向不变集与零超水平集安全集}

控制理论中反复出现一个问题：\textbf{若系统从某个"好的"区域出发，它会一直待在其中吗？} 其数学形式化即正向不变性。考虑自治系统 $\dot x=f(x)$，$f$ 局部 Lipschitz（保证解存在唯一）。

\begin{definition}[正向不变集]\label{def:forward-invariant}
集合 $\mathcal C\subseteq\mathbb R^n$ 称为系统 $\dot x=f(x)$ 的\emph{正向不变集}，若对任意初值 $x(0)\in\mathcal C$，解 $x(t)$ 满足 $x(t)\in\mathcal C$ 对所有 $t\ge0$ 成立。
\end{definition}

直觉是"一旦进入就出不去"——但这是好的陷阱，因为 $\mathcal C$ 正是我们希望系统停留的安全区。在 CBF 理论中，安全集 $\mathcal C$ 通常由标量函数 $h:\mathbb R^n\to\mathbb R$ 的\emph{零超水平集}定义：
\begin{equation}
\mathcal C=\{x:h(x)\ge0\},\quad \partial\mathcal C=\{x:h(x)=0\},\quad \operatorname{Int}(\mathcal C)=\{x:h(x)>0\}.
\label{eq:safe-set}
\end{equation}
要求\textbf{正则性} $\nabla h(x)\neq0$ 对所有 $x\in\partial\mathcal C$，保证边界是光滑超曲面（$n-1$ 维流形）、无尖角奇点。

\begin{insight}[为何用 $h\ge0$ 而非 $h\le c$]
这一符号约定使 $\nabla h$ 在边界上\emph{指向安全集内部}。当 $\dot h=\nabla h\cdot f>0$ 时，速度场与 $\nabla h$ 同向，系统朝内移动。于是后续所有不等式方向统一为"$\ge$"，避免符号混乱——这正是 CBF 理论叙事整洁的根源。
\end{insight}

\begin{example}[两个典型安全集]\label{ex:cbf-safesets}
\emph{圆盘}：机器人位置 $\mathbf p=(x_1,x_2)$ 限于半径 $R$ 的圆盘，取 $h(\mathbf p)=R^2-x_1^2-x_2^2$，则 $h\ge0\iff\|\mathbf p\|\le R$，$\nabla h=(-2x_1,-2x_2)$ 在边界非零。\emph{避障}：障碍在 $\mathbf p_o$、安全距离 $d_{\mathrm{safe}}$，取 $h(\mathbf p)=\|\mathbf p-\mathbf p_o\|^2-d_{\mathrm{safe}}^2$，则 $h\ge0$ 表示到障碍距离不小于 $d_{\mathrm{safe}}$。
\end{example}

\subsection{Bouligand 切锥}

正向不变性的关键在于：\textbf{当状态位于边界时，速度场方向是否指向集合内部（或至少沿切线）？} 若指向外部，系统将于下一时刻离开。

\begin{definition}[Bouligand 切锥]\label{def:tangent-cone}
集合 $\mathcal C$ 在点 $x\in\mathcal C$ 处的切锥为
\[
T_{\mathcal C}(x)=\Big\{v\in\mathbb R^n:\liminf_{\tau\downarrow0}\frac{\mathrm{dist}(x+\tau v,\mathcal C)}{\tau}=0\Big\},\qquad \mathrm{dist}(y,\mathcal C)=\inf_{z\in\mathcal C}\|y-z\|.
\]
\end{definition}

切锥包含所有"不会立刻离开集合"的方向。对水平集 $\mathcal C=\{h\ge0\}$，当 $x\in\partial\mathcal C$（即 $h(x)=0$）且 $\nabla h(x)\neq0$ 时，切锥特化为以 $\nabla h(x)$ 为法向的半空间：
\begin{equation}
T_{\mathcal C}(x)=\{v:\nabla h(x)\cdot v\ge0\},
\label{eq:tangent-levelset}
\end{equation}
恰好是"不让 $h$ 减小"的方向集合。

\subsection{Nagumo 定理}

\begin{theorem}[Nagumo 正不变性，水平集形式 \cite{nagumo1942,blanchini1999survey}]\label{thm:nagumo}
设 $\mathcal C=\{x:h(x)\ge0\}$，$h\in C^1$，$\nabla h(x)\neq0$ 对 $x\in\partial\mathcal C$；系统 $\dot x=f(x)$，$f$ 局部 Lipschitz。则 $\mathcal C$ 前向不变当且仅当
\begin{equation}
\dot h(x)=\nabla h(x)\cdot f(x)\ge0\qquad\forall x\in\partial\mathcal C.
\label{eq:nagumo-cond}
\end{equation}
\end{theorem}

\begin{proof}[必要性]
若存在 $x_0\in\partial\mathcal C$ 使 $\dot h(x_0)=\nabla h(x_0)\cdot f(x_0)<0$，则由 $h(x_0)=0$，对充分小 $\epsilon>0$ 有 $h(x(\epsilon;x_0))=\dot h(x_0)\epsilon+O(\epsilon^2)<0$，轨迹立即离开 $\mathcal C$，与前向不变矛盾。
\end{proof}

\begin{proof}[充分性（思路）]
设条件 \cref{eq:nagumo-cond} 成立。我们采用集合论方法的标准论证：取膨胀的辅助函数证明轨迹无法穿越边界。设 $x(0)\in\mathcal C$，反设存在 $t_1>0$ 使 $h(x(t_1))<0$，令离开时刻 $t^*=\inf\{t>0:h(x(t))<0\}$。由连续性 $h(x(t^*))=0$（即 $x(t^*)\in\partial\mathcal C$）且 $h(x(t))\ge0$ 于 $[0,t^*]$，故右导数 $\dot h(x(t^*))\le0$；又由条件 $\dot h(x(t^*))\ge0$，得 $\dot h(x(t^*))=0$。要排除这种"切线接触后离开"，需对边界邻域作更细的估计：由 $\nabla h\cdot f\ge0$ 在 $\partial\mathcal C$ 上成立且 $f,\nabla h$ 局部 Lipschitz，可构造 $h$ 沿轨迹的下界 ODE（比较引理的变体，见 \cref{ch:lyapunov-basics}），证明 $h(x(t))\ge0$ 在 $t^*$ 的右邻域保持，与 $t^*$ 的定义矛盾。完整论证见 \cite{blanchini1999survey} 对正不变集的集合论刻画。
\end{proof}

\begin{remark}[证明的严谨边界]
上述充分性中"切线接触"情形的闭合依赖比较引理与边界正则性，是教材级证明的精细之处而非数学障碍；Nagumo 定理本身的结论无可置疑。Blanchini 的综述 \cite{blanchini1999survey} 给出了基于切锥与次切导数的完整现代证明，读者若需严格细节可参阅之。
\end{remark}

\subsection{从 Nagumo 到 CBF：推广动机}

\cref{thm:nagumo} 只在边界 $\partial\mathcal C$（$h=0$）上要求 $\dot h\ge0$，对内部（$h>0$）无任何约束。这在\emph{分析}上完全足够（判定一个已知系统是否不变），但在\emph{综合}上不够，原因有二：(1) 构造 QP 约束需要一个\emph{全域}条件，在任意状态 $x$ 处都能判断哪些 $u$ 安全；(2) 仅边界约束意味着系统可以任意高速冲向边界、到边界才"急刹车"，这对有输入约束的系统不现实。

\begin{insight}[反事实推理：只用 Nagumo 会怎样]
若仅用 Nagumo 条件设计控制器，它只在 $h(x)\approx0$ 时才"紧张"并施加修正，$h(x)>0$ 时完全不管安全——犹如一辆车只在即将撞墙时才踩刹车。若车速太快而刹车力有限（输入约束），就会撞上。
\end{insight}

CBF 的推广策略：允许 $h>0$ 时 $\dot h$ 为负（系统可向边界靠近），但衰减速率被 $\alpha(h)$ 限制：
\begin{equation}
\dot h(x)\ge-\alpha(h(x)).
\label{eq:cbf-inequality-intro}
\end{equation}
$h$ 大（远离边界）时 $\alpha(h)$ 大、允许较快靠近；$h$ 小（接近边界）时 $\alpha(h)$ 小、强制减速。下一节将证明比较引理保证此条件下 $h(t)\ge0$。这正像高速公路出口的\emph{递减限速区}：离出口 2\,km 限速 120、1\,km 限速 80、500\,m 限速 60——渐进式约束让车辆即使制动有限也总能安全停下。

\begin{pitfall}[概念误区：混淆 Nagumo 条件与 CBF 条件]
"$\dot h\ge0$ 与 $\dot h\ge-\alpha(h)$ 不就差一个 $\alpha(h)$？"——二者有本质区别。Nagumo 条件 $\dot h\ge0$ 只需在边界满足、内部无要求；CBF 条件 $\dot h\ge-\alpha(h)$ 在\emph{整个状态空间}都有要求（因 $h>0$ 时 $\alpha(h)>0$），从而在远离边界处也限制了接近速率。把安全当成"只在边界处理"的问题是错的：若系统有惯性（如双积分器），到边界才施加最大制动也来不及停下——这正是 \cref{sec:cbf-hocbf} 高阶 CBF 的动机。
\end{pitfall}

% ════════════════════════════════════════════════════════════════
\section{控制 Lyapunov 函数（CLF）}
\label{sec:cbf-clf}

本节把 \cref{ch:lyapunov-basics} 的 Lyapunov 直接法升级为"控制"版本，作为 CBF 的对偶引子。\cref{ch:lyapunov-stability} 给出闭环稳定证书，CLF 则是\emph{开环可镇定}证书：它证明系统\emph{可以}被镇定，但尚未给出具体控制律。

\subsection{CLF 定义与 Artstein 定理}

\begin{definition}[控制 Lyapunov 函数]\label{def:clf}
对控制仿射系统 $\dot x=f(x)+g(x)u$，$C^1$ 函数 $V:\mathbb R^n\to\mathbb R_{\ge0}$ 是 CLF，若它满足 \cref{ch:lyapunov-basics} 中 Lyapunov 函数的基本条件——正定（$V(0)=0$，$V(x)>0$，$x\neq0$）与径向无界（存在 $\alpha_1,\alpha_2\in\mathcal K_\infty$ 使 $\alpha_1(\|x\|)\le V(x)\le\alpha_2(\|x\|)$）——并额外满足\emph{可镇定条件}
\begin{equation}
\inf_{u\in\mathbb R^m}\big[L_fV(x)+L_gV(x)\,u\big]<0\qquad\forall x\neq0.
\label{eq:clf-cond}
\end{equation}
\end{definition}

由于 $L_fV+L_gV\,u$ 关于 $u$ 仿射，上述 $\inf$ 取负当且仅当 $L_gV(x)=0\Rightarrow L_fV(x)<0$，即经典的 \textbf{small control property}：
\begin{equation}
L_gV(x)=0\implies L_fV(x)<0.
\label{eq:scp}
\end{equation}
直觉：若 $L_gV\neq0$，总可沿 $-(L_gV)^\top$ 方向取 $u$ 使 $L_gV\cdot u$ 足够负；唯一需担心的是 $L_gV=0$ 处（控制对 $\dot V$ 无影响），此时只能靠 $f$ 自身让 $V$ 下降，\cref{eq:scp} 恰保证了这一点。

\begin{theorem}[Artstein \cite{artstein1983}]\label{thm:artstein}
对控制仿射系统 $\dot x=f(x)+g(x)u$，以下等价：(i) 系统全局渐近可镇定（存在连续反馈 $u=k(x)$ 使原点全局渐近稳定）；(ii) 存在光滑 CLF $V$。
\end{theorem}

\begin{proof}[思路（双向）]
$\text{(ii)}\Rightarrow\text{(i)}$：给定 CLF，Sontag 公式（\cref{thm:sontag}）显式构造连续反馈 $k(x)$ 使 $\dot V\le-\alpha(V)$。$\text{(i)}\Rightarrow\text{(ii)}$：由逆 Lyapunov 定理（Massera 1956），若闭环 $\dot x=f+gk$ 全局渐近稳定则存在光滑 Lyapunov 函数 $V$；在 $k(x)$ 处 $L_fV+L_gV\,k(x)<0$ 对 $x\neq0$ 成立，而取 $\inf$ 只会更小，故 $V$ 满足 CLF 条件。
\end{proof}

\begin{insight}[Artstein 定理的意义]
它把"能否设计稳定控制器"这一\emph{综合}问题，等价转化为"能否找到满足特定条件的标量函数"这一\emph{分析}问题。一旦找到 CLF，Sontag 公式直接给出控制律。这与 \cref{thm:artstein} 在 CBF 侧的对偶——"找到 CBF $\Rightarrow$ Safety Filter 闭式解"——构成本章的两大构造性支柱。
\end{insight}

\subsection{Sontag 公式：从 CLF 到显式反馈}
\label{sec:sontag}

CLF 只断言"存在 $u$ 使 $\dot V<0$"，未给具体值。最朴素的 $u=-(L_gV)^\top$ 通常过激进且在 $L_gV=0$ 处不连续。Sontag 给出一个优雅的\emph{连续}、满足 small control property 的最小能量解。

设 $a(x)=L_fV(x)$，$b(x)=L_gV(x)\in\mathbb R^{1\times m}$；CLF 条件即 $b=0\Rightarrow a<0$。

\begin{theorem}[Sontag 万能公式 \cite{sontag1989}]\label{thm:sontag}
连续反馈
\begin{equation}
k(x)=\begin{cases}\dfrac{-a(x)-\sqrt{a(x)^2+\|b(x)\|^4}}{\|b(x)\|^2}\,b(x)^\top,& b(x)\neq0,\\[6pt]0,& b(x)=0,\end{cases}
\label{eq:sontag}
\end{equation}
使 $\dot V<0$ 对所有 $x\neq0$，且 $k$ 连续。单输入 $m=1$ 时退化为 $k=\big(-a-\sqrt{a^2+b^4}\big)/b$。
\end{theorem}

\begin{proof}[推导与验证]
\emph{为何取这一形式}：当 $b\neq0$，令 $a+bk=-\sqrt{a^2+\|b\|^4}$。右端恒负（保 $\dot V<0$）；$b\to0$ 且 $a<0$ 时 $u\to0$（连续性）；$a<0$ 时 $|u|$ 小（small control property）。\emph{解出}（单输入）$u=(-a-\sqrt{a^2+b^4})/b$。\emph{验证连续性}：$b\to0$ 时分子分母均趋零，有理化得
\[
u=\frac{-b^3}{-a+\sqrt{a^2+b^4}}\xrightarrow{b\to0}\frac{0}{-a+|a|}=0\quad(\text{因 }a<0,\ -a+|a|=-2a>0).
\]
\emph{验证 $\dot V<0$}（多输入）：取 $k=\frac{-a-\sqrt{a^2+\|b\|^4}}{\|b\|^2}b^\top$，则
\[
\dot V=a+b\,k=a+\frac{-a-\sqrt{a^2+\|b\|^4}}{\|b\|^2}\|b\|^2=-\sqrt{a^2+\|b\|^4}<0,
\]
方向沿 $b^\top$ 是仿射约束下使 $\dot V$ 下降最快的投影；$b=0$ 时取 $u=0$ 仍有 $\dot V=a<0$。
\end{proof}

\begin{pitfall}[数值陷阱：$b\approx0$ 时的除零]
直接用 $u=(-a-\sqrt{a^2+b^4})/b$，当 $|b|<10^{-8}$ 时分母趋零，控制输出可现 $10^{15}$ 量级脉冲、执行器饱和而失稳。\textbf{正确做法}：用有理化等价形式 $u=-b^3/(-a+\sqrt{a^2+b^4})$，或设阈值 $|b|<\epsilon$ 时令 $u=0$。
\end{pitfall}

\subsection{逆最优性与 Sontag\,\texorpdfstring{$\leftrightarrow$}{<->}\,Safety Filter 对偶}
\label{sec:sontag-dual}

Sontag 公式有深刻性质。\textbf{最优性}：它是使 $\dot V=-\sqrt{a^2+\|b\|^4}$ 成立的最小范数 $u$。\textbf{逆最优性}：Freeman 与 Kokotovi\'c \cite{freeman1996inverse} 证明它可视为某隐式最优控制问题 $J=\int_0^\infty[\ell(x)+u^\top R(x)u]\,dt$ 的解，从而赋予其类似 LQR 的鲁棒裕度。最重要的是它与 CBF Safety Filter 的\emph{结构对偶}：

\begin{table}[H]\centering\small
\caption{Sontag 公式与 CBF Safety Filter 的同构结构}
\label{tab:sontag-vs-filter}
\begin{tabular}{lll}
\toprule
& CLF（Sontag） & CBF（Safety Filter，\cref{sec:cbf-qp-filter}） \\
\midrule
目标 & 使 $\dot V$ 足够负 & 使 $\dot h$ 不太负 \\
解的形式 & $u=\lambda\,(L_gV)^\top$ & $u=u_{\mathrm{nom}}+\lambda\,(L_gh)^\top$ \\
$\lambda$ 的含义 & 保收敛的"努力" & 保安全的"修正力度" \\
激活条件 & $L_gV\neq0$ & $a+bu_{\mathrm{nom}}<0$（名义不安全） \\
\bottomrule
\end{tabular}
\end{table}

\begin{insight}[一个对称性，整套框架]
二者在结构上同构：都是沿某梯度方向的投影/修正。Sontag 把 $u$ 从 $0$ 拉到使 $\dot V<0$ 的最小距离；Safety Filter 把 $u_{\mathrm{nom}}$ 投影到使 $\dot h\ge-\alpha(h)$ 的最近点。理解这一对称，就抓住了整个 CLF-CBF 框架的数学本质——\cref{sec:cbf-qp-filter} 将给出 Safety Filter 这一列的闭式解。
\end{insight}

\subsection{CLF-QP}

Sontag 公式不易加输入约束、不易与 CBF 结合、衰减率不可调。现代方法用 QP 替代：
\begin{equation}
u^\star=\arg\min_{u\in\mathcal U}\|u\|^2\quad\text{s.t.}\quad L_fV(x)+L_gV(x)\,u\le-\gamma V(x).
\label{eq:clf-qp}
\end{equation}
取 $\gamma V$ 作衰减率给出指数收敛 $V(t)\le V(0)e^{-\gamma t}$。其可行性条件为 $\exists u\in\mathcal U:L_gV\cdot u\le-L_fV-\gamma V$；当 $L_gV=0$ 时需 $L_fV+\gamma V\le0$（由 \cref{eq:scp} 在 $L_gV=0$ 邻域保证，但 $\gamma$ 过大可能使 $-\gamma V$ 项主导而不可行）。QP 的凸 二次结构、KKT 与求解器细节见 \cref{ch:nlopt}，本章不重述。

% ════════════════════════════════════════════════════════════════
\section{控制屏障函数（CBF）：定义与主定理}
\label{sec:cbf-main}

这是本章第一座理论高峰。我们先把 \cref{ch:lyapunov-basics} 的比较引理特化到安全语境（作为 CBF 的数学引擎），再给出 ZCBF 的严格定义与前向不变性主定理。

\subsection{比较引理的安全特化}
\label{sec:comparison}

CBF 条件 \cref{eq:cbf-inequality-intro} 是一个微分不等式，比较引理把它转化为 $h(t)$ 的显式下界，从而证明 $h$ 不穿零。\cref{ch:lyapunov-basics} 已给出比较引理的一般形式（Khalil 引理 3.4/4.4 \cite{khalil2002nonlinear}）：

\begin{lemma}[比较引理，一般形式，引自 \cref{ch:lyapunov-basics}]\label{lem:comparison}
设 $\dot y(t)\ge\phi(t,y(t))$，$y(t_0)=y_0$，$\phi$ 关于 $y$ 局部 Lipschitz、关于 $t$ 连续；令 $z$ 为 $\dot z=\phi(t,z)$，$z(t_0)=y_0$ 的解。则 $y(t)\ge z(t)$ 对所有解共存的 $t\ge t_0$。对偶版本：$\dot y\le\phi(t,y)\Rightarrow y(t)\le z(t)$。
\end{lemma}

\cref{ch:lyapunov-basics} 用对偶版本从 $\dot V\le-\gamma V$ 推出 $V(t)\le V(0)e^{-\gamma t}$（上界，保收敛）；CBF 用原版从 $\dot h\ge-\alpha(h)$ 推出下界（保不穿零）。同一引理的正反两面，正是 \cref{tab:cbf-clf-duality} 对偶的数学根。取 $\phi(t,y)=-\alpha(y)$，本章只需补证安全特化的关键一环：

\begin{proposition}[比较 ODE 不穿零]\label{prop:no-cross-zero}
设 $\alpha$ 为扩展类 $\mathcal K_\infty$ 函数（连续、严格递增、$\alpha(0)=0$、局部 Lipschitz），$z$ 解 $\dot z=-\alpha(z)$，$z(0)=z_0>0$。则 $z(t)>0$ 对所有有限 $t$。
\end{proposition}

\begin{proof}
$z=0$ 是平衡点（$-\alpha(0)=0$）。由 $\alpha$ 局部 Lipschitz，ODE 解唯一，经过 $(t^*,0)$ 的唯一解是 $z\equiv0$。若从 $z_0>0$ 出发的解在有限时刻 $t^*$ 到达 $z=0$，则它在 $t^*$ 与常值解 $z\equiv0$ 取相同值而此前不同，违反唯一性。故 $z(t)>0$。
\end{proof}

\begin{remark}[非 Lipschitz $\alpha$ 与有限时间收敛]
若 $\alpha$ 在 $0$ 处不 Lipschitz（如 $\alpha(r)=\gamma\sqrt r$），唯一性失效，$z$ 可有限时间到零——这对应 CBF 的\emph{有限时间收敛}变体，在某些应用中有用但需额外分析。线性情形 $\alpha(r)=\gamma r$ 给出 $\dot z=-\gamma z$、显式解 $z(t)=h_0e^{-\gamma t}>0$，故 CBF 保证 $h(t)\ge h_0e^{-\gamma t}$（指数下界），是最常用的选择。
\end{remark}

\subsection{扩展类 \texorpdfstring{$\mathcal K_\infty$}{K-infinity} 函数}

\begin{definition}[扩展类 $\mathcal K_\infty$]\label{def:ext-kinf}
$\alpha:\mathbb R\to\mathbb R$ 是扩展类 $\mathcal K_\infty$ 函数，若 $\alpha$ 连续、严格递增、$\alpha(0)=0$，且 $\alpha(r)\to+\infty$（$r\to+\infty$）、$\alpha(r)\to-\infty$（$r\to-\infty$）。
\end{definition}

与普通类 $\mathcal K_\infty$（仅定义于 $[0,\infty)$）的关键区别是\textbf{定义在整个 $\mathbb R$ 上、含负半轴}。

\begin{insight}[为何必须扩到负半轴]
CBF 条件 $\dot h\ge-\alpha(h)$ 须在 $h<0$（系统暂时不安全）时也有意义。此时 $\alpha(h)<0$，故 $-\alpha(h)>0$，条件变为 $\dot h>0$——\emph{强制 $h$ 增大、把系统拉回安全集}。这正是 ZCBF 名称中 "zeroing"（归零）行为的来源。
\end{insight}

常用选择：线性 $\alpha(r)=\gamma r$（最常用，指数衰减）、立方 $\alpha(r)=\gamma r^3$（零附近更温和）、高阶 $\alpha(r)=\gamma r^{2k-1}$。

\subsection{ZCBF 定义与主定理}

\begin{definition}[Zeroing 控制屏障函数 \cite{ames2017cbf}]\label{def:zcbf}
$h:\mathcal D\to\mathbb R$（$\mathcal D\supseteq\mathcal C$ 开集）是系统的 ZCBF，若存在扩展类 $\mathcal K_\infty$ 函数 $\alpha$ 使
\begin{equation}
\sup_{u\in\mathcal U}\big[L_fh(x)+L_gh(x)\,u\big]\ge-\alpha(h(x))\qquad\forall x\in\mathcal D.
\label{eq:zcbf}
\end{equation}
相应的\emph{可行控制集}为 $K_{\mathrm{cbf}}(x)=\{u\in\mathcal U:L_fh(x)+L_gh(x)\,u\ge-\alpha(h(x))\}$；ZCBF 条件等价于 $K_{\mathrm{cbf}}(x)\neq\varnothing$ 对所有 $x\in\mathcal D$。
\end{definition}

注意 $\sup$ 取在 $u$ 上——不要求\emph{所有} $u$ 满足，只要\emph{存在} $u$ 满足即可。

\begin{theorem}[CBF 前向不变性主定理 \cite{ames2017cbf}]\label{thm:cbf-main}
设 $h$ 为 ZCBF，$\nabla h\neq0$ 于 $\partial\mathcal C$。则任何 Lipschitz 连续控制律 $u(x)\in K_{\mathrm{cbf}}(x)$ 使：(a) $\mathcal C$ 前向不变；(b) $\mathcal C$ 在 $\mathcal D$ 内渐近稳定（$x(0)\notin\mathcal C$ 的轨迹最终回到 $\mathcal C$）。
\end{theorem}

\begin{proof}[Part (a)：前向不变性]
\emph{Step 1（微分不等式）}：取 $u(x)\in K_{\mathrm{cbf}}(x)$ Lipschitz（保闭环 ODE 解存在唯一），沿轨迹 $\dot h(x(t))=L_fh+L_gh\,u\ge-\alpha(h(x(t)))$；令 $y(t)=h(x(t))$，则 $\dot y\ge-\alpha(y)$。\emph{Step 2（比较引理）}：由 \cref{lem:comparison} 取 $\phi=-\alpha$，得 $y(t)\ge z(t)$，其中 $\dot z=-\alpha(z)$，$z(0)=h(x(0))\ge0$。\emph{Step 3（分析比较 ODE）}：若 $z(0)>0$，由 \cref{prop:no-cross-zero} 有 $z(t)>0$；若 $z(0)=0$，由唯一性 $z\equiv0$。两情形均 $z(t)\ge0$，故
\[
h(x(t))=y(t)\ge z(t)\ge0\qquad\forall t\ge0,
\]
即 $x(t)\in\mathcal C$ 保持。
\end{proof}

\begin{proof}[Part (b)：zeroing 渐近稳定]
若 $x(0)\notin\mathcal C$，即 $z(0)=h(x(0))<0$。由 $\alpha(z)<0$ 当 $z<0$，得 $\dot z=-\alpha(z)>0$，$z$ 递增趋向平衡点 $0$（标准 Lipschitz $\alpha$ 下不有限时间到达）。故 $y(t)=h(x(t))\ge z(t)\to0^-$，轨迹被拉回安全集边界。
\end{proof}

\begin{insight}[一枚硬币的两面]
CBF 条件的几何含义：把 $h$ 想成"安全裕度仪表"，CBF 说"仪表下降速度不超过比较 ODE 给出的衰减速度"；而比较 ODE 从正初值出发永远为正，故仪表永远为正、系统永远安全。这与 CLF 条件 $\dot V\le-\gamma(V)$ 精确对偶：CLF 说"Lyapunov 仪表至少以 $\gamma(V)$ 下降"（保收敛），CBF 说"安全仪表至多以 $\alpha(h)$ 下降"（保不穿零）。
\end{insight}

\subsection{\texorpdfstring{$\alpha$}{alpha} 函数的选择与工程影响}

$\alpha$ 的选择直接塑造控制器行为。

\begin{table}[H]\centering\small
\caption{典型 $\alpha$ 函数及其比较 ODE 解与行为}
\label{tab:alpha-choice}
\begin{tabular}{llll}
\toprule
$\alpha(h)$ & CBF 约束 & 比较 ODE 解 & 行为 \\
\midrule
$\gamma h$（线性） & $\dot h\ge-\gamma h$ & $h(t)\ge h_0e^{-\gamma t}$ & 指数衰减，最常用 \\
$\gamma h^3$（立方） & $\dot h\ge-\gamma h^3$ & 更慢衰减 & 远离边界时更宽松 \\
$\gamma\sqrt h$（根号，$h\ge0$） & $\dot h\ge-\gamma\sqrt h$ & 有限时间到零 & 更保守，$h=0$ 处非 Lipschitz \\
\bottomrule
\end{tabular}
\end{table}

\textbf{$\gamma$ 的工程意义}：$\gamma$ 大（如 $10$）允许 $h$ 快速衰减、系统可高速接近边界，CBF 对名义控制侵入小，但接近边界时需急刹车、输入有限则可能来不及；$\gamma$ 小（如 $0.5$）限制 $h$ 缓慢衰减、从远处即减速，安全余量大但过度保守。\textbf{时间常数} $\tau=1/\gamma$ 决定从安全态到达边界附近的最坏情况时间：若系统有制动延迟 $\tau_d$，应选 $\gamma$ 使 $\tau\gg\tau_d$，否则在安全裕度耗尽前来不及响应。工程上多从 $\alpha(r)=\gamma r$、$\gamma\in[1,10]$ 起调；高频安全层（约 $0.1$–$1$\,kHz）配 $\gamma=O(1)$ 量级通常给出充裕响应时间。也可用复合形式 $\alpha(h)=\gamma_1 h+\gamma_2 h^3$（远处立方项主导、近处线性项温和）或 $\alpha(h)=\gamma\tanh(h/\epsilon)$（把约束饱和在 $[-\gamma,\gamma]$，避免 $h$ 很大时约束形同虚设）。

\subsection{Reciprocal CBF：为何 ZCBF 成事实标准}

早期的\emph{倒数型屏障函数}（Reciprocal CBF, RCBF）取 $B:\operatorname{Int}(\mathcal C)\to\mathbb R_{\ge0}$，$B\to+\infty$ 当 $x\to\partial\mathcal C$（典型 $B=1/h$），条件 $\inf_u[L_fB+L_gB\,u]\le\alpha(1/B)$。

\begin{table}[H]\centering\small
\caption{Zeroing CBF 与 Reciprocal CBF 对比}
\label{tab:zcbf-vs-rcbf}
\begin{tabular}{lll}
\toprule
特性 & Zeroing CBF & Reciprocal CBF \\
\midrule
值域 & $h:\mathbb R^n\to\mathbb R$（有限） & $B:\operatorname{Int}(\mathcal C)\to[0,\infty)$（边界爆破） \\
边界行为 & $h\to0$（光滑） & $B\to\infty$（发散） \\
梯度 & 边界有界（$\nabla h\neq0$ 但有限） & 边界发散（$\nabla B\to\infty$） \\
数值性质 & 友好（QP 系数有界） & 不友好（梯度爆炸，QP 病态） \\
不安全域 & 自然定义（$h<0$ 有意义） & 无定义（$B$ 仅在 $\operatorname{Int}$） \\
现代使用 & 事实标准 & 已基本弃用 \\
\bottomrule
\end{tabular}
\end{table}

\begin{insight}[反事实推理：若用 RCBF]
系统接近边界时 $B=1/h\to\infty$、$\nabla B=-\nabla h/h^2\to\infty$，QP 约束系数 $L_gB$ 爆炸，导致数值不稳定、解出的 $u$ 大而方向不合理、执行器饱和后安全性丧失。ZCBF 的 $\nabla h$ 在边界因正则性保持有界，QP 始终良态——这是 ZCBF 取代 RCBF 成为标准的根本原因。本章其余部分一律采用 ZCBF。
\end{insight}

% ════════════════════════════════════════════════════════════════
\section{高相对度 CBF：HOCBF 与 ECBF}
\label{sec:cbf-hocbf}

第二座理论高峰。当控制不直接出现在 $\dot h$ 中（相对度 $>1$），标准 CBF 失效，须用高阶构造。

\subsection{标准 CBF 的失效与相对度}

考虑双积分器 $\ddot q=u$（即 $\dot x_1=x_2$，$\dot x_2=u$，$x_1=q$、$x_2=\dot q$），位置约束 $h=x_1-q_{\min}$。计算 $L_fh=[1,0]\cdot[x_2,0]^\top=x_2$，$L_gh=[1,0]\cdot[0,1]^\top=0$。\textbf{灾难}：$L_gh\equiv0$，控制 $u$ 根本不出现在 $\dot h$ 中，CBF 约束 $x_2\ge-\alpha(h)$ 是对\emph{状态}的约束、无法写进 QP。物理上：位置约束的安全性取决于 $q$，而 $u$ 直接影响的是加速度 $\ddot q$，从加速度到位置需两次积分——这就是"相对度 2"。

\begin{definition}[相对度]\label{def:rel-degree}
安全约束 $h$ 相对于 $\dot x=f(x)+g(x)u$ 的\emph{相对度} $r$ 是使 $L_gL_f^{r-1}h(x)\neq0$ 的最小正整数，即 $u$ 首次出现在 $h^{(r)}=L_f^rh+L_gL_f^{r-1}h\cdot u$ 中。对 $r>1$，$L_gh=L_gL_fh=\cdots=L_gL_f^{r-2}h=0$，标准 CBF 失效。
\end{definition}

\begin{pitfall}[陷阱：未检查相对度就用标准 CBF]
对双积分器位置约束直接写 $L_gh\cdot u\ge-\alpha h$，则 $L_gh$ 恒为 $0$，约束退化成 $0\ge-\alpha h$——当 $h>0$ 时恒成立，等于没有约束，系统直接撞上。\textbf{检测方法}：算 $L_gh$，若为零向量则相对度 $\ge2$，须改用 HOCBF/ECBF。
\end{pitfall}

\subsection{HOCBF 递归构造}

\textbf{核心思想}：不直接约束 $h$，而是构造一条"约束链" $\psi_0,\dots,\psi_{r-1}$，每层约束上一层的导数，最终在第 $r$ 层引入 $u$。

\begin{definition}[高阶 CBF \cite{xiao2022hocbf}]\label{def:hocbf}
令 $\psi_0(x):=h(x)$，并递归
\begin{equation}
\psi_i(x):=\dot\psi_{i-1}(x)+\alpha_i(\psi_{i-1}(x)),\qquad i=1,\dots,r-1,
\label{eq:hocbf-recur}
\end{equation}
其中 $\alpha_i$ 为扩展类 $\mathcal K_\infty$ 函数（可任意非线性）。HOCBF 约束只在第 $r$ 层施加（此时 $u$ 出现）：
\begin{equation}
\dot\psi_{r-1}(x,u)+\alpha_r(\psi_{r-1}(x))\ge0,
\label{eq:hocbf-constraint}
\end{equation}
展开即 $L_f^rh+L_gL_f^{r-1}h\cdot u+(\text{含 }h,L_fh,\dots,\alpha_i\text{ 的项})\ge0$，是关于 $u$ 的线性约束。
\end{definition}

\textbf{安全性保证}：若控制器满足 \cref{eq:hocbf-constraint}，则所有 $\mathcal C_i=\{\psi_{i-1}\ge0\}$ 前向不变；由 $\psi_0=h$，$\mathcal C_0=\{h\ge0\}$ 前向不变，原始安全集得保。

\begin{insight}[嵌套的安全缓冲区]
HOCBF 像嵌套的安全缓冲。设想一辆驶向悬崖的车：第 $0$ 层（$h\ge0$）车不能掉下悬崖（位置）；第 $1$ 层（$\psi_1\ge0$）即使还没到崖边，速度太快也不行（速度约束，取决于到崖的距离）；第 $2$ 层（约束本身）制动力必须足够大以保证减速（加速度）。越远离悬崖允许速度越大，越接近要求制动力越大——层层限速直至控制层。
\end{insight}

\subsection{双积分器完整算例}

\begin{example}[双积分器位置限制的 HOCBF]\label{ex:hocbf-double}
系统 $\dot x_1=x_2$，$\dot x_2=u$，约束 $h=x_1-q_{\min}$，$q_{\min}=-1$。取线性 $\alpha_i(r)=\gamma_i r$、$\gamma_1=\gamma_2=2$：
\[
\psi_0=x_1+1,\qquad \psi_1=\dot\psi_0+\gamma_1\psi_0=x_2+2(x_1+1)=x_2+2x_1+2.
\]
$\psi_1\ge0$ 意味着"当前速度加上基于位置裕度的补偿非负"：即使 $x_2<0$（向下），只要位置裕度 $x_1+1\ge|x_2|/2$ 仍算安全速度。HOCBF 约束 $\dot\psi_1+\gamma_2\psi_1\ge0$，其中 $\dot\psi_1=u+2x_2$，故
\[
u+2x_2+2(x_2+2x_1+2)\ge0\ \Longrightarrow\ u\ge-4(x_1+x_2+1).
\]
\end{example}

\begin{table}[H]\centering\small
\caption{\cref{ex:hocbf-double} 中 $u$ 下界随状态的变化}
\label{tab:hocbf-numerical}
\begin{tabular}{llll}
\toprule
状态 $(x_1,x_2)$ & 物理含义 & $u$ 下界 & 解读 \\
\midrule
$(0,-1)$ & 位置在 $0$、向下 & $0$ & 必须开始减速 \\
$(-0.5,-2)$ & 接近下限、快速下降 & $6$ & 紧急制动 \\
$(5,0)$ & 远离边界、静止 & $-24$ & 几乎不约束 \\
$(-0.9,-0.1)$ & 极近边界、慢速 & $0$ & 恰处临界 \\
\bottomrule
\end{tabular}
\end{table}

几何上，$\mathcal C_0=\{x_1\ge-1\}$ 与 $\mathcal C_1=\{x_2\ge-2x_1-2\}$ 均为半平面，可行域 $\mathcal C_0\cap\mathcal C_1$ 在相平面上是楔形区域：一旦系统在交集内（且 HOCBF 约束被满足）就永远在其中。\textbf{与极点配置的联系}：$\gamma_1=\gamma_2=2$ 对应特征方程 $(s+2)^2=0$、双重极点 $s=-2$；约束激活时系统行为如同被极点 $-2$ 约束的二阶系统。

\begin{pitfall}[概念误区：$\psi_1$ 的符号]
"$\psi_1=\dot h+\alpha_1(h)$，当 $\dot h<0$ 但 $|\dot h|<\alpha_1(h)$ 时 $\psi_1>0$——为何速度为负也算安全？"——$\psi_1\ge0$ 不是要求 $\dot h\ge0$（那是 Nagumo），而是 $\dot h\ge-\alpha_1(h)$：允许 $h$ 下降（向边界接近），但速率被 $\alpha_1(h)$ 约束。这正是 CBF 相对 Nagumo 的推广在每一层的体现。
\end{pitfall}

\subsection{指数 CBF（ECBF）}
\label{sec:ecbf}

指数 CBF 是 Nguyen 与 Sreenath \cite{nguyen2016ecbf} 提出的方法，是 HOCBF 当所有 $\alpha_i$ 取线性时的特例，其优势是可用极点配置设计参数——这与 \cref{ch:linear-synthesis} 的极点配置/可控标准形直接对接。

对相对度 $r$ 的约束 $h$，定义状态向量 $\eta_b(x)=[h,L_fh,\dots,L_f^{r-1}h]^\top\in\mathbb R^r$。动力学化为 Brunovsky 标准形（积分器链）$\dot\eta_b=F\eta_b+G\mu$，其中 $F$ 是 $r\times r$ 移位矩阵、$G=[0,\dots,0,1]^\top$、$\mu=L_f^rh+L_gL_f^{r-1}h\cdot u$。ECBF 约束为
\begin{equation}
L_f^rh+L_gL_f^{r-1}h\cdot u\ge-K_b\,\eta_b(x),
\label{eq:ecbf}
\end{equation}
其中 $K_b=[k_1,\dots,k_r]\in\mathbb R^{1\times r}$ 通过\textbf{极点配置}选择，使 $F-GK_b$ 的特征值全在左半平面（Hurwitz）。以相对度 $2$ 为例，特征多项式 $s^2+k_2s+k_1=0$，选重极点 $s=-p$ 得 $k_1=p^2$、$k_2=2p$，约束化为 $u\ge-p^2h-2p\dot h$（当 $L_gL_fh>0$ 时除以它）：$p$ 大则回复快但控制量大、$p$ 小则温和。\textbf{ECBF 与 HOCBF 的关系}：当 HOCBF 中 $\alpha_i(r)=k_i r$ 全取线性时二者等价；HOCBF 允许非线性 $\alpha_i$（如 $\sqrt\cdot$、$\arctan$），提供更大设计灵活性。

% ════════════════════════════════════════════════════════════════
\section{CLF-CBF-QP 统一框架}
\label{sec:cbf-qp}

至此我们有了 CLF（稳定证书）与 CBF（安全证书），本节把二者熔进一个实时可解的 QP，并给出其 KKT 闭式解与最常用的 Safety Filter 变体。

\subsection{冲突分析与优先级不对称}

我们希望控制器同时\emph{渐近稳定}（CLF）且\emph{永远安全}（CBF）。二者可能冲突：机器人需到达目标但路径上有障碍——CLF 指引它直冲目标（穿障），CBF 要求避障。冲突的数学表现是 $K_{\mathrm{clf}}(x)\cap K_{\mathrm{cbf}}(x)=\varnothing$，不存在 $u$ 同时满足 CLF 约束 $L_fV+L_gV\,u\le-\gamma V$ 与 CBF 约束 $L_fh+L_gh\,u\ge-\alpha h$。

\begin{insight}[安全优先的数学实现]
机器人学中安全与稳定的优先级\emph{不对称}：不安全可致灾难（撞毁、伤人），不稳定只是性能下降（收敛变慢）。因此 CBF 作硬约束（不可违反），CLF 经松弛变量降为软约束。这是"安全优先"原则的数学落地。
\end{insight}

\subsection{CLF-CBF-QP 的构造}

\begin{equation}
\begin{aligned}
(u^\star,\delta^\star)=\arg\min_{u\in\mathbb R^m,\,\delta\in\mathbb R}\quad&\tfrac12\|u-u_{\mathrm{ref}}\|_H^2+p\,\delta^2\\
\text{s.t.}\quad&L_fV+L_gV\,u\le-\gamma V+\delta &&(\text{CLF, 软})\\
&L_fh+L_gh\,u\ge-\alpha(h) &&(\text{CBF, 硬})\\
&u_{\min}\le u\le u_{\max} &&(\text{输入约束})
\end{aligned}
\label{eq:clf-cbf-qp}
\end{equation}
其中 $H\succ0$ 是控制代价权重，$u_{\mathrm{ref}}$ 为参考控制（PD 输出/MPC 解/RL 策略），$p\in[10^2,10^4]$ 为松弛惩罚权，$\delta$ 为 CLF 松弛变量，$\gamma$ 为 CLF 衰减率。

\begin{remark}[为何用 $p\delta^2$ 而非 $p|\delta|$]
二次惩罚保持目标二次（凸 QP 结构）且自动给出 $\delta=0$ 的软趋势；线性惩罚 $p|\delta|$ 会使问题变成二次锥规划（SOCP），求解更复杂。
\end{remark}

\subsection{KKT 条件与四情形闭式解}

为深入理解 QP 行为，简化为无输入约束情形（$\mathcal U=\mathbb R^m$，$H=I$，$u_{\mathrm{ref}}=0$）。拉格朗日量
\[
\mathcal L=\tfrac12\|u\|^2+p\delta^2+\lambda_V(L_fV+L_gV\,u+\gamma V-\delta)+\lambda_h(-L_fh-L_gh\,u-\alpha(h)).
\]
KKT 条件（一般理论见 \cref{ch:nlopt}）给出：驻点 $u^\star=-\lambda_V(L_gV)^\top+\lambda_h(L_gh)^\top$、$\delta^\star=\lambda_V/(2p)$；对偶可行 $\lambda_V,\lambda_h\ge0$；以及原始可行与互补松弛 $\lambda_V(\cdots)=0$、$\lambda_h(\cdots)=0$。按乘子激活分四情形：

\begin{table}[H]\centering\small
\caption{CLF-CBF-QP 的 KKT 四情形}
\label{tab:kkt-cases}
\begin{tabular}{lllll}
\toprule
情形 & $\lambda_V$ & $\lambda_h$ & 物理含义 & $u^\star$ \\
\midrule
两约束均不激活 & $0$ & $0$ & 名义控制已安全且收敛 & $u_{\mathrm{ref}}$ \\
仅 CLF 激活 & $>0$ & $0$ & 远离障碍、需修正保收敛 & $u_{\mathrm{ref}}-\lambda_V(L_gV)^\top$ \\
仅 CBF 激活 & $0$ & $>0$ & 已收敛但接近危险 & $u_{\mathrm{ref}}+\lambda_h(L_gh)^\top$ \\
两者均激活 & $>0$ & $>0$ & 冲突！牺牲稳定保安全 & 两力平衡 \\
\bottomrule
\end{tabular}
\end{table}

冲突时 $\delta^\star=\lambda_V/(2p)>0$，CLF 约束被软化（允许 $\dot V>-\gamma V$）；$p$ 越大则 $\delta$ 越小、CLF 越接近硬约束，但可能导致 QP 不可行（硬 CLF + 硬 CBF 不相容）。

\begin{pitfall}[松弛罚权 $p$ 设置不当]
$p$ 太小（如 $1$）：$\delta$ 可很大、CLF 几乎失效，系统只保安全不收敛（机器人停在原地）。$p$ 太大（如 $10^8$）：$\delta\approx0$、CLF 近似硬约束，与 CBF 冲突时 QP 不可行。\textbf{推荐} $p\in[10^2,10^4]$；实践中先设大 $p$，频繁不可行则减小。
\end{pitfall}

\subsection{Safety Filter 变体与闭式解}
\label{sec:cbf-qp-filter}

工程中最常用的形式：给定名义控制 $u_{\mathrm{nom}}(x)$（来自 RL 策略、MPC、遥操作、PD），做\emph{最小修正}保安全：
\begin{equation}
u^\star=\arg\min_{u\in\mathcal U}\|u-u_{\mathrm{nom}}\|^2\quad\text{s.t.}\quad L_fh+L_gh\,u\ge-\alpha(h).
\label{eq:safety-filter}
\end{equation}
\textbf{最小侵入原则}：只在必须时修正、修正量最小化；RL 策略在安全时完全不受干扰，仅在即将违规时被轻推回安全侧。

\begin{theorem}[Safety Filter 的闭式解]\label{thm:filter-closed}
设 $\mathcal U=\mathbb R^m$，记 $a=L_fh(x)+\alpha(h(x))$，$b=L_gh(x)\in\mathbb R^{1\times m}$，CBF 约束即 $a+bu\ge0$。则
\begin{equation}
u^\star=u_{\mathrm{nom}}+\lambda^\star b^\top,\qquad \lambda^\star=\max\Big\{0,\ \frac{-(a+bu_{\mathrm{nom}})}{bb^\top}\Big\}.
\label{eq:filter-closed}
\end{equation}
即名义安全（$a+bu_{\mathrm{nom}}\ge0$）时 $u^\star=u_{\mathrm{nom}}$ 不修正；否则 $u^\star$ 是 $u_{\mathrm{nom}}$ 在超平面 $\{u:a+bu=0\}$ 上的正交投影。
\end{theorem}

\begin{proof}[几何推导]
\cref{eq:safety-filter} 是把 $u_{\mathrm{nom}}$ 投影到半空间 $\{u:a+bu\ge0\}$。若 $u_{\mathrm{nom}}$ 已在该半空间内，投影即自身。否则最近点落在边界超平面上，投影方向沿其法向 $b^\top$，写 $u^\star=u_{\mathrm{nom}}+\lambda b^\top$ 并令 $a+bu^\star=0$ 解得 $\lambda=-(a+bu_{\mathrm{nom}})/(bb^\top)$；$\max\{0,\cdot\}$ 合并两情形。
\end{proof}

这个闭式解计算量极小：一次内积（检查安全性）加一次向量加法（修正），无输入约束时甚至无需 QP 求解器。\textbf{Lipschitz 连续性}：当 $bb^\top>0$（即 $L_gh\neq0$，相对度 $1$ 的自然要求），$u^\star$ 关于 $x$ Lipschitz 连续，保证闭环 ODE 解存在唯一 \cite{morris2013lipschitz}。回看 \cref{tab:sontag-vs-filter}：\cref{eq:filter-closed} 与 Sontag 公式 \cref{eq:sontag} 结构同构，正是 §\ref{sec:sontag-dual} 所许诺的对偶。

\begin{pitfall}[误区：Safety Filter "改变了"策略]
Safety Filter 只在\emph{必要时}修正，大部分时间 $u^\star=u_{\mathrm{nom}}$；若 RL 策略本身已安全，它永不激活，不改变学习目标或训练过程。\textbf{关键区别}：Safety Filter 是\emph{部署时}（deployment-time）的后处理，\textbf{不是}训练时约束；训练时的安全约束是 Constrained RL（如 CPO/PCPO），二者互补（详见 \cref{sec:cbf-unify}）。
\end{pitfall}

\subsection{多障碍物与 Composite CBF}

若有 $N$ 个安全约束 $h_1,\dots,h_N$，直接在 QP 中加 $N$ 条 $L_fh_i+L_gh_i\,u\ge-\alpha_i(h_i)$，保证 $\mathcal C=\bigcap_i\{h_i\ge0\}$ 前向不变；可行性条件为 $\bigcap_iK_{\mathrm{cbf}}^i(x)\cap\mathcal U\neq\varnothing$。当约束过多或安全集过小时可能不可行。

为避免 QP 规模随障碍数线性增长，可把多个 $h_i$ \emph{合并}为单一约束。多 CBF 的非光滑 $\min/\max$（布尔）合并由 Glotfelter、Cort\'es 与 Egerstedt \cite{glotfelter2017nonsmooth} 首倡；其\emph{光滑化}（log-sum-exp）版本——composite CBF——可写为
\begin{equation}
h(x)=-\frac1\kappa\log\sum_i e^{-\kappa h_i(x)}\ \xrightarrow{\ \kappa\to\infty\ }\ \min_i h_i(x),
\label{eq:composite-cbf}
\end{equation}
合并后仅一条 CBF 约束（QP 规模不随 $N$ 增），代价是 $\nabla h$ 涉及所有 $h_i$ 梯度的加权和、可能较保守 \cite{molnar2024composing}。

\subsection{QP 到求解器标准形的映射}

实时实现把 \cref{eq:clf-cbf-qp} 映射到求解器（如 OSQP）的标准形 $\min_\chi\tfrac12\chi^\top P\chi+q^\top\chi$ s.t. $l\le A\chi\le u$。决策变量 $\chi=[u^\top,\delta]^\top\in\mathbb R^{m+1}$，
\begin{equation}
P=\begin{bmatrix}H&0\\0&2p\end{bmatrix},\quad q=\begin{bmatrix}-H u_{\mathrm{ref}}\\0\end{bmatrix},\quad
A=\begin{bmatrix}L_gV&-1\\-L_gh&0\\I_m&0\end{bmatrix},\ l=\begin{bmatrix}-\infty\\-\infty\\u_{\min}\end{bmatrix},\ u=\begin{bmatrix}-L_fV-\gamma V\\L_fh+\alpha(h)\\u_{\max}\end{bmatrix}.
\label{eq:qp-cbf-standard}
\end{equation}
ADMM/active-set 等具体凸 QP 求解、warm-start 与代码生成属 \cref{ch:nlopt} 范畴，本章不展开。典型规模下（单输入到十余维）QP 求解在微秒到毫秒级，足以支撑 $0.1$–$10$\,kHz 的安全层。

% ════════════════════════════════════════════════════════════════
\section{可行性与控制不变集}
\label{sec:cbf-feasibility}

CBF 理论中最实际的挑战：当输入有界时，是否总存在安全控制？本节把答案接到 \cref{ch:hj-reachability} 的可达性理论。

\subsection{输入约束相容性}

\begin{definition}[控制不变集]\label{def:control-invariant}
$\mathcal C$ 是控制不变的，若对所有 $x\in\mathcal C$，存在 $u\in\mathcal U$ 使 $x$ 留在 $\mathcal C$ 中，即 $K_{\mathrm{cbf}}(x)\cap\mathcal U\neq\varnothing$ 对所有 $x\in\mathcal C$。
\end{definition}

对 $\mathcal U=\{u:\|u\|\le u_{\max}\}$，$\max_{u\in\mathcal U}L_gh\cdot u=\|L_gh\|\,u_{\max}$，故可行性条件为
\begin{equation}
L_fh(x)+\|L_gh(x)\|\,u_{\max}\ge-\alpha(h(x)).
\label{eq:input-feasibility}
\end{equation}

\subsection{可行性丧失与最大控制不变集}

\cref{eq:input-feasibility} 在某状态不满足，说明该状态"不可救"。物理场景：(1) \emph{制动力不足}——车以 100\,km/h 冲向 5\,m 外墙壁，即使最大制动也无法避免；(2) \emph{多约束夹击}——机器人被三面墙包围，任何方向都违反某 CBF；(3) \emph{高速接近}——当前 $h$ 仍为正但所需制动超出能力。

\begin{insight}[viability kernel：CBF 与 HJ 的交汇]
使内部所有状态都可行的\emph{最大}集合，正是 \cref{ch:hj-reachability} 的\emph{最大控制不变集}（viability kernel）。CBF 在此集内保证可行；HJ 可达性通过求解 HJI-VI 精确计算该集（但维度受限）。二者关系：CBF 是廉价的充分条件、HJ 是昂贵的充要刻画。工程上应从更高层（路径规划）避免进入 viability kernel 之外的区域——这是预防 QP 不可行的根本手段。HJI-VI 与粘性解的主体在 \cref{ch:hj-reachability}，本章仅借其结论。
\end{insight}

\subsection{未期望平衡点与死锁}

\begin{theorem}[CLF-CBF-QP 引入未期望平衡点 \cite{reis2021undesired}]\label{thm:undesired-eq}
CLF-CBF-QP 闭环会在安全集边界上引入除 CLF 极小点外的\emph{渐近稳定平衡点}。
\end{theorem}

物理现象：机器人卡在障碍物前无法绕行——CLF 指向目标（穿障）、CBF 阻止靠近，两力平衡处成为"陷阱"。\textbf{解法}：引入旋转场（rotation-based guidance）打破对称；由高层路径规划提供绕行 $u_{\mathrm{ref}}$；用导航函数（navigation function）设计无局部极小的 CLF。这是 CBF 理论的重要局限，提示安全滤波须与全局规划协同——规划部对此有专门处理，本章留前向钩子。

% ════════════════════════════════════════════════════════════════
\section{离散与鲁棒 CBF}
\label{sec:cbf-robust}

前几节假设连续时间、精确模型。真实机器人是采样控制且模型有误差，本节补上离散化与不确定性两块拼图。

\subsection{离散时间 CBF（DCBF）}

真实控制器以固定频率（如 1\,kHz）运行，两次更新间系统在零阶保持下演化；连续 CBF 条件在离散实现中只近似成立——采样间隔 $\Delta t$ 大或动态快时，连续 CBF 不能保证离散系统安全（如 $\dot x=10u$、$\Delta t=0.1$\,s，$t\in(0,0.1)$ 内可能已越界）。

\begin{definition}[离散时间 CBF \cite{agrawal2017dcbf}]\label{def:dcbf}
对离散系统 $x_{k+1}=F(x_k,u_k)$、安全集 $\mathcal C=\{h\ge0\}$，$h$ 是 DCBF，若存在 $\alpha\in(0,1]$ 使 $\sup_{u_k\in\mathcal U}h(F(x_k,u_k))\ge(1-\alpha)h(x_k)$，等价约束 $h(x_{k+1})\ge(1-\alpha)h(x_k)$。
\end{definition}

$\alpha$ 的三档解读：$\alpha=0$ 要求 $h(x_{k+1})\ge h(x_k)$（裕度永不减，过保守）；$\alpha=1$ 要求 $h(x_{k+1})\ge0$（下一步安全即可，最宽松）；$\alpha\in(0,1)$ 给几何衰减 $h(x_k)\ge(1-\alpha)^k h(x_0)$（指数收敛到零但不穿零）。\textbf{与连续 CBF 的关系}：$\dot h\ge-\gamma h$ 的 Euler 离散化为 $h_{k+1}\ge(1-\gamma\Delta t)h_k$，对应 $\alpha=\gamma\Delta t$；当 $\gamma\Delta t>1$ 时失去意义——这正是连续 CBF 在低频采样下失效的原因。

\begin{remark}[DCBF 一般是 NLP 而非 QP]
约束 $h(F(x_k,u_k))\ge(1-\alpha)h(x_k)$ 关于 $u_k$ 通常\emph{非线性}（因 $F$ 非线性），不能直接写成 QP。处理：一阶 Taylor 近似 $h(F)\approx h(F(x_k,0))+\nabla_uh(F)\cdot u_k$ 回到线性约束；或在 MPC 框架内作 NLP 约束（即 §\ref{sec:cbf-unify} 的 MPC-CBF）；若 $F(x,u)=F_0(x)+G(x)u$ 仿射，则 $h(F)$ 可能仍是 $u$ 的仿射函数。
\end{remark}

\subsection{输入到状态安全（ISSf-CBF）}

真实系统有模型误差 $\dot x=f(x)+g(x)u+d(t)$，$\|d\|\le d_{\max}$。\textbf{输入到状态安全}（ISSf）是 \cref{ch:lyapunov-stability} 中 ISS 的安全对偶：ISS 谓 $\dot V\le-\alpha(\|x\|)+\sigma(\|u\|)$ 意味"扰动有界则状态有界"；ISSf 把同样思想翻译到安全语言。

\begin{theorem}[ISSf-CBF \cite{kolathaya2019issf}]\label{thm:issf}
若 $\sup_u[L_fh+L_gh\,u]\ge-\alpha(h(x))+\iota(\|d\|)$（$\iota\in\mathcal K$），则扰动下 $\mathcal C$ 的一个\emph{膨胀集} $\mathcal C_d\supseteq\mathcal C$ 前向不变，膨胀量随 $d_{\max}$ 增大。
\end{theorem}

直觉：扰动使安全边界"模糊化"，实际安全集比标称的小。工程实现是在 CBF 约束中加 margin：
\begin{equation}
L_fh+L_gh\,u\ge-\alpha(h)+d_{\max}\|\nabla h\|,
\label{eq:issf-margin}
\end{equation}
右端多出的 $d_{\max}\|\nabla h\|$ 补偿最坏情况扰动对 $\dot h$ 的影响。ISS 的定义与一般理论在 \cref{ch:lyapunov-stability}，本章仅给其安全特化。

\subsection{时变 CBF}

当障碍移动或环境变化时，$h$ 本身随时间变化为 $h(x,t)$，CBF 条件变为
\begin{equation}
\frac{\partial h}{\partial t}+L_fh+L_gh\,u\ge-\alpha(h).
\label{eq:time-varying-cbf}
\end{equation}
额外项 $\partial h/\partial t$ 表示环境变化对安全裕度的影响。对已知运动模型的障碍 $\dot{\mathbf p}_o=\mathbf v_o$，$\partial h/\partial t=-2(\mathbf p-\mathbf p_o)^\top\mathbf v_o$；预测不确定时用可达集或最坏情况；感知到控制的延迟 $\tau$ 内障碍可能移动 $\|\mathbf v_o\|\tau$，需额外 margin。

\begin{insight}[反事实推理：忽略 $\partial h/\partial t$]
对动态障碍只用静态 CBF，则障碍快速靠近（$\partial h/\partial t\ll0$）时，实际 $\dot h$ 比控制器估计的更负——系统自以为安全而实际不安全。这是自动驾驶中 CBF 用于行人避障的核心挑战。
\end{insight}

% ════════════════════════════════════════════════════════════════
\section{与 HJ / MPC / 学习的统一}
\label{sec:cbf-unify}

前几节假设模型精确、$h$ 可手工设计。真实系统模型有误差、环境有噪声、高维 $h$ 无法手工构造。本节勾勒突破这些限制的前沿，并把 CBF 接到本书其他章。

\subsection{CBVF：CBF 与 HJ 可达的桥}

HJ 可达性通过求解 HJI-VI 获得\emph{最大鲁棒安全集}（viability kernel），给出最优解但计算量随维度指数增长；CBF 计算廉价（只需 QP）但需手工设计 $h$。\textbf{控制屏障-值函数}（CBVF, Choi 等 \cite{choi2021cbvf}）统一二者：引入折扣 $\gamma_{\mathrm{cbvf}}\ge0$，CBVF $B_{\gamma_{\mathrm{cbvf}}}(x,t)$ 是如下 HJI-VI 的粘性解：
\begin{equation}
0=\min\Big\{\ell(x)-B_{\gamma_{\mathrm{cbvf}}},\ \partial_tB_{\gamma_{\mathrm{cbvf}}}+\max_{u\in\mathcal U}\min_{d\in\mathcal D}\nabla_xB_{\gamma_{\mathrm{cbvf}}}\cdot f(x,u,d)+\gamma_{\mathrm{cbvf}}B_{\gamma_{\mathrm{cbvf}}}\Big\}.
\label{eq:cbvf}
\end{equation}
关键事实：$\gamma_{\mathrm{cbvf}}=0$ 退化为经典 HJ 可达，$\{B_0\ge0\}$ 是 viability kernel；$\gamma_{\mathrm{cbvf}}>0$ 内嵌 CBF 衰减 $\alpha(r)=\gamma_{\mathrm{cbvf}}r$，$\{B_{\gamma_{\mathrm{cbvf}}}\ge0\}$ 仍精确等于 viability kernel。实践上离线用 HJ 方法算 $B_{\gamma_{\mathrm{cbvf}}}$、在线用鲁棒 CBVF-QP，解决了手工设计 CBF 的痛点，缺点是 HJ 只能处理低维（实用上限约 $n\le6$）。HJI-VI 与粘性解的主证在 \cref{ch:hj-reachability}（其 value-function 视角亦与 \cref{ch:dp-hjb} 的 HJB 同源），本节仅给结论。

\subsection{MPC-CBF}

CBF-QP 只看当前一步、无预见性；MPC 有有限时域预测但计算重。Zeng、Zhang 与 Sreenath \cite{zeng2021mpccbf} 把 DCBF 作为 MPC 每一预测步的硬约束：
\[
\min_{u_0,\dots,u_{N-1}}\sum_{k=0}^{N-1}\ell(x_k,u_k)+V_f(x_N)\quad\text{s.t.}\quad x_{k+1}=F(x_k,u_k),\ h(x_{k+1})\ge(1-\alpha)h(x_k),\ u_k\in\mathcal U.
\]
优势：短预测时域（$N=5\sim10$）即可实现远距离避障，因 DCBF 每步保证裕度不快速衰减、即便没预测到很远也安全。MPC 的递归可行性、tube/robust 外层理论由 \cref{ch:mpc-advanced} 承接，本章不重述。

\subsection{CBF 作 RL Safety Filter 与可微 QP 层}

RL 通过试错探索，"错"在真机上可能意味碰撞、坠落、损坏。CBF-QP 作 Safety Filter 是当前最成熟方案 \cite{cheng2019rlcbf}：策略 $\pi_\theta(s)$ 输出 $u_{\mathrm{nom}}$，经 CBF-QP 滤为 $u_{\mathrm{safe}}$ 再施于环境。\textbf{可微 QP 层}是关键技巧：用 \texttt{cvxpylayers} \cite{agrawal2019cvxpylayers} 或 \texttt{qpth} \cite{amos2017optnet} 把 QP 封装为可微层，反向传播通过 KKT 条件的隐函数微分，使策略梯度
\[
\nabla_\theta J=\mathbb E\Big[\nabla_\theta\log\pi_\theta(u_{\mathrm{nom}})\cdot\frac{\partial u_{\mathrm{safe}}}{\partial u_{\mathrm{nom}}}\cdot R\Big]
\]
中的 $\partial u_{\mathrm{safe}}/\partial u_{\mathrm{nom}}$ 经 QP 的 KKT 隐式微分获得（KKT 隐微分见 \cref{ch:nlopt}）。工程效果：策略在训练中学会"预判"安全约束、逐渐减少 Filter 激活——即学会安全行为而非依赖后处理。须区分：\textbf{部署时}（Safety Filter，本节）与\textbf{训练时}（Constrained RL，如 CPO/PCPO）是两条互补的线；强化学习运控部展开后者，本章只给前者的理论一段。

\subsection{Neural CBF 与高维验证}

对高维系统（如 12-DoF 四旋翼、视觉感知移动机器人），手工设计 $h$ 几乎不可能。\textbf{Neural CBF} 把 $h_\theta$ 参数化为神经网络。Dawson、Qin、Gao 与 Fan \cite{dawson2021safe} 给出三项训练损失：
\begin{equation}
\mathcal L(\theta)=\underbrace{\lambda_1\mathcal L_{\mathrm{class}}}_{\text{安全/不安全分类}}+\underbrace{\lambda_2\mathcal L_{\mathrm{descent}}}_{\text{CBF 条件满足}}+\underbrace{\lambda_3\mathcal L_{\mathrm{boundary}}}_{\text{边界正则}},
\label{eq:neural-cbf-loss}
\end{equation}
其中 $\mathcal L_{\mathrm{class}}$ 要求 $h_\theta>0$ 于安全演示、$<0$ 于碰撞状态；$\mathcal L_{\mathrm{descent}}$ 是 $\max(0,-L_fh_\theta-L_gh_\theta u-\alpha(h_\theta))$ 的采样均值；$\mathcal L_{\mathrm{boundary}}$ 鼓励 $\{h_\theta=0\}$ 接近真实安全边界。训练后用 SMT（dReal/Z3）或 SOS 验证 $h_\theta$ 在状态空间满足 CBF 条件，失败则用反例点增广训练集重训。神经证书的统一视角综述见 \cite{dawson2023survey}。

\begin{table}[H]\centering\small
\caption{Neural CBF 验证方法的维度谱系}
\label{tab:neural-cbf-verify}
\begin{tabular}{llll}
\toprule
方法 & 适用维度 & 保证类型 & 代表工作 \\
\midrule
SMT（dReal/Z3） & $n\le4$ & 形式化完备 & \cite{dawson2021safe} \\
SOS（多项式系统） & $n\le6$ & 形式化 & \cite{prajna2007barrier} \\
Lipschitz 边界 & $n\le10$ & 概率 + Lipschitz & \cite{fazlyab2019lipschitz} \\
随机采样（PAC bound） & 任意 $n$ & 概率 & \cite{robey2020learning} \\
反例引导学习 & 任意 $n$ & 启发式 & --- \\
\bottomrule
\end{tabular}
\end{table}

\subsection{随机 CBF 与数据驱动 CBF}

对随机系统（传感器噪声、风场无人机），确定性安全不可能、只能保\emph{概率安全}。对 Ito SDE $dx=(f+gu)\,dt+\sigma(x)\,dW_t$，Ito 公式对屏障 $B$ 求广义导数（与确定性的关键区别在多出扩散项）：
\begin{equation}
dB(x_t)=\Big[\nabla B\cdot(f+gu)+\underbrace{\tfrac12\operatorname{tr}(\sigma^\top\nabla^2B\,\sigma)}_{\text{扩散项（Ito 修正）}}\Big]dt+\nabla B\cdot\sigma\,dW_t.
\label{eq:ito-barrier}
\end{equation}
扩散项 $\tfrac12\operatorname{tr}(\sigma^\top\nabla^2B\,\sigma)$ 是噪声对 $B$ 的"平均推动力"：即使漂移方向安全，噪声的二阶效应（经 $B$ 的曲率 $\nabla^2B$）也可能把系统推向不安全——这是随机系统比确定性更"危险"的数学原因。若随机 CBF 条件 $\sup_u[\nabla B\cdot(f+gu)+\tfrac12\operatorname{tr}(\sigma^\top\nabla^2B\,\sigma)]\ge-\alpha(B)$ 成立，则 $B(x_t)$ 是\emph{上鞅}（supermartingale），由 Doob 最大不等式得概率界
\begin{equation}
\mathbb P\Big(\inf_{t\ge0}h(x_t)<0\Big)\le\frac{B(x_0)}{c}.
\label{eq:doob-bound}
\end{equation}

\begin{remark}[随机安全须用上鞅概率界，勿声称"概率 1"]
\cref{eq:doob-bound} 用的是倒数型 $B$ 的上鞅 + Doob 不等式给出的\emph{概率}界（有效），适用于工程上预留"噪声裕度带"。但须警惕：关于 zeroing 型随机 CBF "几乎必然（概率 1）安全"的若干结论已被指出存在缺陷，故本书\textbf{不}把随机 ZCBF 写成"概率 1 安全"——稳妥的保证是 \cref{eq:doob-bound} 这样的上鞅概率界（Prajna 等 \cite{prajna2007barrier}、Clark \cite{clark2021stochastic}）。
\end{remark}

当模型 $f,g$ 不精确时，$L_fh$、$L_gh$ 计算有误差。\textbf{GP-CBF}（Khojasteh 等 \cite{khojasteh2020gpcbf}）用高斯过程估计动力学残差 $\Delta f\sim\mathcal{GP}(0,k)$，后验给出均值 $\mu$ 与方差 $\sigma^2$，约束变为 $L_{\hat f}h+L_gh\,u+\mu^\top\nabla h-\beta\sigma\|\nabla h\|\ge-\alpha(h)$（$\beta$ 控制置信度，$\beta=2$ 约 95\%），方差大的区域约束更紧（不确定时更谨慎）。\textbf{从演示学 CBF}（Robey 等 \cite{robey2020learning}）则从专家安全/不安全轨迹训练分类器 $h_\theta$，在 Lipschitz 假设下给可证明安全保证。

% ════════════════════════════════════════════════════════════════
\section{应用实例}
\label{sec:cbf-apps}

理论之后，三个 worked example 把 CBF 从公式落到具体系统：ACC（相对度 1 入门）、四旋翼着陆（HOCBF 相对度 2）、多机器人避障。

\subsection{自适应巡航控制（ACC）}
\label{sec:cbf-acc}

ACC 是理解 CBF 最直觉的入口，也是 CBF-QP 的旗舰算例——首见于 Ames、Grizzle 与 Tabuada 2014 年的工作 \cite{ames2014cbfacc}，后被 2017 年的奠基论文 \cite{ames2017cbf} 沿用。

\begin{example}[ACC 的 CBF-QP]\label{ex:cbf-acc}
高速公路两车跟驰，状态 $x=[z,v_e,v_l]^\top$：$z=p_l-p_e$ 为车距、$v_e$ 自车速度、$v_l$ 前车速度。动力学 $\dot z=v_l-v_e$、$\dot v_e=u$、$\dot v_l=a_l(t)$（前车加速度，不可控），即 $f=[v_l-v_e,0,a_l]^\top$、$g=[0,1,0]^\top$。\emph{Time-headway 准则}\footnote{ISO 15622:2018 规定 ACC 跟车时距下限 $\tau_{\min}\ge0.8$\,s，并要求在车速高于 $8$\,m/s 时至少提供一档落在 $1.5$–$2.2$\,s 区间的设定；本例取 $\tau_h=1.8$\,s 作该区间内的代表值，并非标准强制的唯一时距。}（取时距 $\tau_h=1.8$\,s）给安全约束
\[
h(x)=z-\tau_h v_e=(p_l-p_e)-1.8\,v_e.
\]
\emph{相对度}：$L_gh=[1,-1.8,0]\cdot[0,1,0]^\top=-1.8\neq0$，相对度 $1$。选 $\alpha(h)=\gamma h$，CBF 约束展开：
\[
(v_l-v_e)-1.8\,u\ge-\gamma(z-1.8v_e)\ \Longrightarrow\ u\le\frac{(v_l-v_e)+\gamma(z-1.8v_e)}{1.8}.
\]
\emph{数值}：取 $v_e=30$、$v_l=25$、$z=60$（m/s, m）、$\gamma=1$，则裕度 $h=60-1.8\times30=6$\,m，最大允许加速度 $u_{\max}=(25-30+6)/1.8=0.56$\,m/s$^2$。若名义 $u_{\mathrm{nom}}=2$\,m/s$^2$，Safety Filter 修正为 $u^\star=0.56$\,m/s$^2$。
\end{example}

\begin{note}[教学简化说明]
\cref{ex:cbf-acc} 取 $\dot v_e=u$（控制仿射 $g=[0,1,0]^\top$）略去了滚动阻力，Ames 等的完整模型含 $\dot v_e=-\tfrac1mF_r(v_e)+\tfrac1mu$（$F_r=f_0+f_1v+f_2v^2$）\cite{ames2014cbfacc}。这是合理的教学简化，物理解读不变：当 $v_l<v_e$（前车更慢）且车距不富裕时，最大允许加速度可为负（要求减速）；$\gamma$ 控制"多早开始减速"——$\gamma$ 大则更早减速更保守。
\end{note}

\subsection{二维单积分器避障}

\begin{example}[单积分器避圆障]\label{ex:cbf-avoid}
系统 $\dot{\mathbf p}=u$，$\mathbf p=[x,y]^\top$（$f=0$，$g=I$）；障碍圆心 $\mathbf p_o=[3,3]$、半径 $R=1$。取 $h(\mathbf p)=\|\mathbf p-\mathbf p_o\|^2-R^2=(x-3)^2+(y-3)^2-1$，则 $\nabla h=[2(x-3),2(y-3)]$、$L_fh=0$、$L_gh=\nabla h$，相对度 $1$。CBF 约束（$\alpha(h)=h$）为 $2(x-3)u_x+2(y-3)u_y\ge-h$，Safety Filter $\min\|u-u_{\mathrm{nom}}\|^2$ s.t. 此线性约束，闭式解由 \cref{thm:filter-closed} 给出。多障碍则叠加多条约束或用 \cref{eq:composite-cbf} 合并。
\end{example}

\subsection{四旋翼安全着陆（HOCBF）}

\begin{example}[四旋翼高度 HOCBF]\label{ex:cbf-quad}
四旋翼简化为质心平移（忽略姿态环、设内环快）$\ddot{\mathbf p}=\tfrac1m F_{\mathrm{thrust}}-\mathsf g\,e_3+d$，状态 $x=[\mathbf p^\top,\dot{\mathbf p}^\top]^\top\in\mathbb R^6$、控制 $u=F_{\mathrm{thrust}}\in\mathbb R^3$（此处重力记 $\mathsf g$ 以区别输入矩阵 $g(x)$）。高度下限 $h_1=z-z_{\min}$：$\dot h_1=\dot z$（不含 $u$，$L_gh_1=0$）、$\ddot h_1=F_z/m-\mathsf g$（含 $u$），\emph{相对度 2}，须用 HOCBF。构造 $\psi_0=z-z_{\min}$、$\psi_1=\dot z+\gamma_1(z-z_{\min})$，约束 $\dot\psi_1+\gamma_2\psi_1\ge0$ 展开为
\[
\frac{F_z}{m}-\mathsf g+(\gamma_1+\gamma_2)\dot z+\gamma_1\gamma_2(z-z_{\min})\ge0\ \Longrightarrow\ F_z\ge m\big[\mathsf g-(\gamma_1+\gamma_2)\dot z-\gamma_1\gamma_2(z-z_{\min})\big].
\]
极点配置取 $\gamma_1=\gamma_2=2$（重极点 $s=-2$）得 $F_z\ge m[\mathsf g-4\dot z-4(z-z_{\min})]$：当 $z$ 接近 $z_{\min}$ 且 $\dot z<0$（下降中）时约束要求 $F_z$ 增大，即加推力减缓下降。
\end{example}

\subsection{多机器人避障}

\begin{example}[成对 CBF 多机避障]\label{ex:cbf-multi}
为 $N$ 个单积分器机器人两两定义 $h_{ij}=\|\mathbf p_i-\mathbf p_j\|^2-(2R)^2$（$2R$ 为最小允许间距）\cite{wang2017safety}。\emph{集中式}：所有 $N(N-1)/2$ 约束放入一个 $Nm$ 维 QP，保证全局安全但计算量 $O(N^3)$。\emph{分布式}：每机器人只考虑邻居，但机器人 $i,j$ 对同一约束各自优化，可能都想左闪而相撞；解法是 consistent perturbation——规定优先级（如 ID 小者优先），低优先级者在其约束中额外考虑高优先级者的预期动作。
\end{example}

\textbf{死锁与对称性破缺}：两机器人正面相遇时对称 CBF 给对称修正、都停下而死锁。解法：非对称优先级（ID 小者让路方向固定）、加小随机偏置打破对称、或高层规划避免正面相遇。多机器人安全的更完整处理见规划部，本章给出通用 CBF 内核，留前向钩子。

% ════════════════════════════════════════════════════════════════
\section*{本章常见误解汇总}

\begin{table}[H]\centering\small
\caption{CLF/CBF/QP 安全控制常见误解与正解}
\label{tab:cbf-misconceptions}
\begin{tabular}{p{0.46\linewidth}p{0.46\linewidth}}
\toprule
\textbf{常见误解} & \textbf{正确理解} \\
\midrule
$\dot h\ge0$（Nagumo）与 $\dot h\ge-\alpha(h)$（CBF）只差一项 & 前者仅边界约束、后者全域约束；后者在远离边界处也限制接近速率 \\
相对度未检查就用标准 CBF & $L_gh\equiv0$ 时约束失效（$0\ge-\alpha h$ 恒成立），须用 HOCBF/ECBF \\
Safety Filter 改变了 RL 策略 & 只在必要时最小修正，是部署时后处理，非训练时约束 \\
随机 ZCBF 保证"概率 1"安全 & 稳妥保证是上鞅 + Doob 概率界（\cref{eq:doob-bound}）；a.s. 结论存争议 \\
CLF-CBF-QP 一定收敛到目标 & 边界可能引入未期望平衡点/死锁（\cref{thm:undesired-eq}），须配全局规划 \\
RCBF 与 ZCBF 等价可换 & RCBF 边界梯度爆炸、QP 病态；ZCBF 边界有界，是事实标准 \\
\bottomrule
\end{tabular}
\end{table}

% ════════════════════════════════════════════════════════════════
\section*{本章小结}

本章把"安全"建立为与稳定性同等严格的控制证书，主线为：集合不变性（Nagumo）$\to$ CLF 引子 $\to$ CBF 主体 $\to$ 高相对度 $\to$ 统一 QP $\to$ 可行性 $\to$ 离散/鲁棒 $\to$ 前沿统一 $\to$ 应用。

\begin{table}[H]\centering\small
\caption{本章核心知识点速查}
\label{tab:cbf-summary}
\begin{tabular}{lll}
\toprule
知识点 & 核心公式/条件 & 关键直觉 \\
\midrule
正向不变集（\cref{def:forward-invariant}） & $x(0)\in\mathcal C\Rightarrow x(t)\in\mathcal C$ & 进去就出不来 \\
Nagumo（\cref{thm:nagumo}） & $\dot h\ge0$ 于 $\partial\mathcal C$ & 边界处速度不指向外 \\
CLF（\cref{def:clf}） & $\inf_u[L_fV+L_gVu]<0$ & 存在控制使 $V$ 下降 \\
Sontag（\cref{thm:sontag}） & $k=\frac{-a-\sqrt{a^2+\|b\|^4}}{\|b\|^2}b^\top$ & CLF 的显式连续反馈 \\
ZCBF（\cref{def:zcbf}） & $\sup_u[L_fh+L_ghu]\ge-\alpha(h)$ & 存在控制使 $h$ 不快降 \\
CBF 主定理（\cref{thm:cbf-main}） & 比较引理 $\Rightarrow h(t)\ge0$ & 衰减率不超比较 ODE \\
HOCBF（\cref{def:hocbf}） & 递归 $\psi_i$，约束在第 $r$ 层 & 层层限速到控制层 \\
CLF-CBF-QP（\cref{eq:clf-cbf-qp}） & 软 CLF + 硬 CBF & 安全优先于稳定 \\
Safety Filter（\cref{thm:filter-closed}） & $u^\star=u_{\mathrm{nom}}+\lambda^\star b^\top$ & 最小侵入修正 \\
可行性（\cref{eq:input-feasibility}） & $L_fh+\|L_gh\|u_{\max}\ge-\alpha(h)$ & 制动力是否够 \\
ISSf-CBF（\cref{thm:issf}） & 加 margin $d_{\max}\|\nabla h\|$ & 安全集膨胀应对扰动 \\
CBVF（\cref{eq:cbvf}） & 折扣 HJI-VI，$\{B\ge0\}$=viability kernel & CBF 与 HJ 的桥 \\
\bottomrule
\end{tabular}
\end{table}

三个核心证书的不等式并置，最能体现本章的对偶美感：CLF 条件 $L_fV+L_gV\,u\le-\gamma V$（保收敛）、CBF 条件 $L_fh+L_gh\,u\ge-\alpha(h)$（保安全）、CLF-CBF-QP 把二者熔为软硬两类约束的实时优化。理解 \cref{tab:cbf-clf-duality} 的逐项对偶与 \cref{tab:sontag-vs-filter} 的结构同构，便掌握了整个框架的数学本质。

% ════════════════════════════════════════════════════════════════
\section*{通向高级：本章交给后续什么}

本章把 CBF 安全证书做成了通用内核，它向四个方向延伸：(1) 向 \cref{ch:hj-reachability}——CBVF（\cref{eq:cbvf}）与可行性的 viability kernel（\cref{def:control-invariant}）都以 HJI-VI 为精确刻画，本章给桥、I 章给主证；(2) 向 \cref{ch:mpc-advanced}——MPC-CBF 把 DCBF 嵌入预测时域，本章给思想、G 章承接递归可行性与 tube/robust 安全滤波；(3) 向规划部——CBF 安全滤波在不确定性规划中作场景应用，本章给理论、规划部给应用，互引去重，且其微分博弈/HJI 与本章 CBVF 同源；(4) 向强化学习运控部——CBF-QP 作部署时安全层 + 可微 KKT 层是 safe-RL 主流范式，本章给部署时一段、RL 部给训练时全景。此外，CBF 约束可嵌入力控/WBC 的分层 QP（与操作空间、力位混合控制的 QP 同源），为接触力安全提供统一接口。至此，从"会收敛"到"永远安全"的安全控制理论已自成体系，可作为后续一切安全关键控制的地基。

% ════════════════════════════════════════════════════════════════
\section*{延伸阅读}

\begin{itemize}
\item \textbf{入门与综述}：Ames 等的 ECC 2019 教程 \cite{ames2019cbf} 以约 40 页覆盖 CBF 全貌，是最佳入门；Khalil《Nonlinear Systems》\cite{khalil2002nonlinear} 第 4 章给 Lyapunov/比较引理的严格基础；Blanchini 的正不变性综述 \cite{blanchini1999survey} 是集合不变性理论的权威。
\item \textbf{核心论文}：ZCBF/RCBF 统一定义与 CLF-CBF-QP 主定理见 \cite{ames2017cbf}；HOCBF 递归构造见 \cite{xiao2022hocbf}；ECBF 见 \cite{nguyen2016ecbf}；CBVF 见 \cite{choi2021cbvf}；MPC-CBF 见 \cite{zeng2021mpccbf}；多机器人 CBF 见 \cite{wang2017safety}。
\item \textbf{学习与安全 RL}：CBF 作 RL Safety Filter 见 \cite{cheng2019rlcbf}；Neural CBF 三项损失见 \cite{dawson2021safe}、其统一视角综述见 \cite{dawson2023survey}；safe-RL 全景综述见 \cite{brunke2022safe}。
\item \textbf{开源实现}：CBF-QP 可由 OSQP（QP 求解）、cbfkit/cbfpy（CBF-QP，含 JAX）、hj\_reachability（HJ PDE）、neural\_clbf（神经证书）等库支撑，工程读者可据此快速搭建安全层原型。
\end{itemize}
